\subsection*{Three falsifiable predictions}
\label{sec:disc:predictions}

The model generates three empirically testable predictions that are currently unverified in any bee
species:

\textbf{Prediction 1 (anti-correlation)}: Correlating GCA battery scores~\cite{AnimalCog2024} with
opt-out accuracy in the same individuals should yield a negative correlation ($r < -0.3$). The
test requires same-individual measurements across at least two paradigms: a multi-task GCA battery
and an opt-out uncertainty monitoring task~\cite{Giurfa2013}. Both paradigms exist; what is missing
is their application to the same marked individuals.

\textbf{Prediction 2 (inverted-U dose-response)}: Testing opt-out accuracy across a range of
precision levels (induced by varying training intensity, reward certainty, or mild pharmacological
challenge) should reveal a non-monotonic relationship with a maximum at intermediate precision.
A multi-dose neonicotinoid protocol~\cite{SciReports2020} could modulate precision, with the
prediction that mild doses improve opt-out accuracy while severe doses impair it.

\textbf{Prediction 3 (disruption sign flip)}: Under circadian disruption (constant light or
neonicotinoid exposure), the metacognition--GCA anti-correlation should reverse sign. This requires
running both the GCA battery and the opt-out task on disrupted and control animals and comparing
the cross-individual correlation structure.

\subsection*{Relation to the Dunning-Kruger effect}
\label{sec:disc:dunning}

The overconfidence paradox bears surface similarity to the Dunning-Kruger
effect~\cite{DunningKruger1999}, but the mechanism is different. The Dunning-Kruger effect
(in humans) proposes that incompetent individuals overestimate their ability due to lack of
metacognitive skills. Our model predicts that highly \emph{competent} agents underperform on
metacognitive tasks because precision suppresses the uncertainty signal. This is an inverted
Dunning-Kruger: it is the skilled, not the unskilled, who have degraded self-knowledge. Empirical
tests in high-precision-trained bees vs. naive bees could distinguish these accounts.

\subsection*{Limitations}
\label{sec:disc:limits}

Several limitations should be noted. First, precision is modeled as a scalar; the biological CX
almost certainly employs multiple, partially independent precision channels (spatial vs. temporal
vs. olfactory). A multi-channel extension would require different predictions for domain-specific
disruptions. Second, we model behavior at the individual level without a neural implementation;
connecting to actual CX topology (ring attractor, Kenyon cells) is left for future work. Third,
the model has no learning dynamics: precision is drawn at the start of each simulation and does
not change within trials. Including online precision updating would make the model richer but
introduces additional parameters. Finally, we emphasize that all model parameters are grounded in
published data~\cite{AvarWeber2012}, but empirical validation requires measuring the same
individuals across all four behavioral domains, which has not yet been done.

\subsection*{Relation to the Free Energy Principle}
\label{sec:disc:fep}

Our model is a behavioral-level abstraction of the FEP~\cite{Friston2010}. Under FEP, precision
plays the role of attention: agents allocate precision to minimize expected free energy. Our
precision parameter captures this role at the population level, without implementing the
variational inference that FEP specifies at the neural level. A natural extension would map our
scalar precision to the FEP's sensory precision hyperparameter, deriving predictions about
neuromodulatory control of the CX (e.g., octopamine vs. dopamine as precision-adjusting
neuromodulators in bees~\cite{Perry2013}).

\subsection*{Broader implications}
\label{sec:disc:broader}

If the predictions are confirmed, this model has implications beyond bees. The precision--metacognition
trade-off may be a general property of any system that uses a common precision signal across
performance and self-monitoring. In artificial agents, an analogous tension between predictive
accuracy and uncertainty quantification has been noted in Bayesian deep learning~\cite{Clark2013}.
In humans, individuals with autism spectrum disorder show high local precision but impaired global
uncertainty monitoring---a pattern consistent with the model's predictions at the high-precision
extreme.
